\documentclass{article}
\usepackage{amsmath,amssymb,amsthm,hyperref,url}
\newtheorem{theorem}{Theorem}
\newtheorem{proposition}[theorem]{Proposition}
\newtheorem{lemma}[theorem]{Lemma}
\DeclareMathOperator{\CVaR}{CVaR}
\DeclareMathOperator{\osc}{osc}
\DeclareMathOperator{\col}{col}
\title{Corrected quadratic payments for a fixed CVaR surrogate}
\author{}
\date{}
\begin{document}
\maketitle
\begin{abstract}
% Ledger: 1--9, 13--19, 20, 22, 24, 26.
We examine the combination of a quadratic resource-allocation mechanism
and an empirical conditional value-at-risk (CVaR) oracle. The mechanism's
strict matrix certificate is incompatible with its implementation equalities,
and its printed payment can violate individual rationality. An adjustment
depending only on opponents' messages repairs equilibrium transfers.
Under additional curvature and multiplier-existence assumptions, a weighted
monotonicity argument implements the unique optimum of a fixed expected
empirical CVaR surrogate, with exact budget balance and surrogate individual
rationality. The oscillation of the surrogate error bounds population-risk
regret, participation shortfall and welfare loss. Contracting only the loss
queries preserves the zero outside action and local feasibility. A standard
nested stochastic proximal-point construction gives mean-square allocation
convergence to this fixed surrogate optimum, without requiring unique prices.
These conclusions neither validate the original best-response learner nor
establish exact population-risk learning.
\end{abstract}

\section{The question and its scope}
% Ledger: 3, 12, 15, 17, 19, 24, 26.
Garjani, Shokri and Kebriaei~\cite{Q}, called Q below, propose quadratic
payments for strategic allocation with capacity constraints. Wang, Huang,
Xu and Johansson~\cite{P}, called P, develop a cooperative distributed
CVaR method using loss values. P establishes consensus and convergence to a
fixed expected empirical smoothed objective under its fixed-batch assumptions.
Q's implementation and dynamic guarantees are conditional on certificates
and assumptions that cannot simply be carried into that setting.

% Ledger: 1--9, 13, 17, 19, 24, 26.
Our contribution is a correction and conditional combination of these
constructions. We give a direct implementation proof for Q's explicit
payment after an incentive-preserving transfer adjustment, and quantify
how P's statistical surrogate changes economic guarantees. The new content
is the specific repair and its combination with the economic error bounds.
The value estimator, elementary smoothing construction and generic
proximal-point convergence principle are not claimed as new. Throughout,
implementation means Nash implementation in the message game; no
truthful-reporting or dominant-strategy assertion is made.

\section{Payments and implementation}
\subsection{Model and the printed payment}
% Ledger: 4, 6--9, 17, 19; source: Q, Section III, label eq40.
There are \(N\ge2\) agents and \(r\) resources. Agent \(n\) chooses
\(s_n=(x_n,p_n)\in X_n\times\mathbb R_+^r\), where \(X_n\) is compact
and convex. Capacity is \(c\in\mathbb R_+^r\). The planner minimises
\(\sum_n f_n(x_n)\) over
\[
 \mathcal A=\{x\in\textstyle\prod_n X_n:\ \sum_nx_n\le c\}.
\]
The normal cone convention is
\(N_X(x)=\{v:\langle v,y-x\rangle\le0\text{ for all }y\in X\}\).
A planner Karush--Kuhn--Tucker (KKT) pair \((x^\star,\lambda)\) satisfies
\begin{equation}\label{eq:kkt}
 0\in\partial f_n(x_n^\star)+\lambda+N_{X_n}(x_n^\star),\quad
 \lambda\ge0,\quad x^\star\in\mathcal A,\quad
 \lambda^\top(\sum_nx_n^\star-c)=0.
\end{equation}
The game cost is \(f_n(x_n)+t_n(s)\); the coupling constraint is imposed
at equilibrium, not on unilateral message choices.

% Ledger: 4--9, 17; source: Q, label eq40, reproduced to specify the game.
Write \(\gamma>0\) for Q's payment scale, and define
\(P_{-n}=\sum_{j\ne n}p_j\), \(X_{-n}=\sum_{j\ne n}x_j\),
\(S_{-n}=\sum_{j\ne n}\|p_j\|^2\), and
\(D_{-n}=\sum_{j\ne n}p_j^\top x_j\).
The printed payment in Q, Section III, equation labelled
\texttt{eq40} in its source, is
\begin{equation}\label{eq:tax}
\begin{aligned}
 t_n^Q(s)=\gamma\bigg[&\frac12\|p_n\|^2
 -\frac{p_n^\top P_{-n}}{N-1}
 -\frac{p_n^\top(x_n+X_{-n}-c)}{N-1}\\
 &+\frac{N P_{-n}^\top(x_n-c/N)}{(N-1)^2}
 +\frac{S_{-n}}{2(N-1)}
 -\frac{D_{-n}-P_{-n}^\top c/N}{(N-1)^2}\bigg].
\end{aligned}
\end{equation}
All derivatives below are taken from this expression.

\subsection{Obstructions in Q}
% Ledger: 1--3, 5--7, 9.
\begin{proposition}\label{prop:failure}
Q's strict certificate P1(ii) and implementation equality P2(i) cannot
hold simultaneously. Moreover, payment~\eqref{eq:tax} need not be
individually rational, even with normalised strongly concave valuations.
\end{proposition}
\begin{proof}
% Ledger: 1--3.
In Q's notation, P2(i) requires \(\sum_m A_{nm}^n=0\).
For \(v\ne0\) and \(h=\col((0,v),\ldots,(0,v))\), its price blocks give
\[
 h^\top(\Psi+\Psi^\top)h
 =-2\sum_n v^\top\Big(\sum_m A_{nm}^n\Big)v=0.
\]
P1(ii), \(\Psi+\Psi^\top\preceq-\upsilon I\) with \(\upsilon>0\),
would instead bound this by \(-\upsilon N\|v\|^2<0\).
This disproves simultaneous feasibility of those sufficient conditions.

% Ledger: 6, 7, 9.
For the payment claim, take \(N=3\), \(\gamma=c=1\), \(X_n=[0,1]\),
and valuations
\[
 V_1(x)=10x-x^2/2,\qquad V_2(x)=V_3(x)=0.1x-x^2/2,
 \qquad f_n=-V_n.
\]
At \(x=(1,0,0)\), \(p=(9,9,9)\), own price derivatives vanish and
own allocation cost derivatives are \(0,8.9,8.9\). The positive derivatives
occur at lower bounds. Each own cost Hessian is
\(\left(\begin{smallmatrix}1&-1/2\\-1/2&1\end{smallmatrix}\right)\succ0\),
so these conditions suffice for a Nash equilibrium. Substitution gives
\(t_1^Q=10.5>V_1(1)=9.5\).
\end{proof}
% Ledger: 5--9.
More generally, put \(b_N=(N^2-N+1)/(N-1)^2\). At common prices
\(p_n=p\) satisfying \(p^\top(\sum_nx_n-c)=0\), collection of terms gives
\begin{equation}\label{eq:old-transfer}
 t_n^Q=\gamma b_Np^\top(x_n-c/N).
\end{equation}
The printed linear coefficients also satisfy
\(\sum_m a_m^n=\gamma c/[N(N-1)]\), violating Q's P4(ii) when
capacity has a positive component. These failures do not imply absence
of equilibria or impossibility of every implementing mechanism.

\subsection{A direct replacement theorem}
% Ledger: 4, 8, 9, 14, 17, 19.
Assume each \(f_n\) is continuous and \(m\)-strongly convex on \(X_n\),
with convex extensions sufficient for the subgradient and normal-cone sum
rules. Assume a planner KKT pair exists. Define
\begin{equation}\label{eq:repair}
 t_n^{\rm corr}=t_n^Q+
 \gamma(b_N-1)\left(\frac{P_{-n}}{N-1}\right)^\top
 \left(X_{-n}-\frac{N-1}{N}c\right).
\end{equation}
% Ledger: 8, 9, 17, 19.
\begin{theorem}\label{thm:implementation}
If
\begin{equation}\label{eq:threshold}
 m>\frac{\gamma(2N-1)^2}{4(N-1)^4},
\end{equation}
then every equilibrium of the corrected game implements the unique planner
allocation. Prices are common and \(\lambda=\gamma p_n\) is a planner
multiplier. Conversely, every planner KKT pair gives an equilibrium at
\(p_n=\lambda/\gamma\). If also \(0\in X_n\) and \(f_n(0)=0\),
then for \(V_n=-f_n\) the corrected equilibrium payments satisfy
\begin{equation}\label{eq:transfers}
 t_n^{\rm corr}=\lambda^\top(x_n-c/N),\qquad
 \sum_nt_n^{\rm corr}=0,\qquad
 V_n(x_n)-t_n^{\rm corr}\ge\lambda^\top c/N\ge0.
\end{equation}
Prices need not be unique.
\end{theorem}
\begin{proof}
% Ledger: 4, 9, 14, 17, 19.
The added term depends only on opponents, preserving every best response.
Stack allocations and prices separately. Let
\(\Pi=\mathbf1\mathbf1^\top/N\), \(\Delta=I-\Pi\), and
\(\kappa=(2N-1)/(N-1)^2\); these matrices act resourcewise.
Differentiating~\eqref{eq:tax} and multiplying price rows by \((N-1)/N\)
gives the weighted own-cost operator
\begin{equation}\label{eq:operator}
 G_x\in\partial f(x)+\gamma(\Pi-\kappa\Delta)p,\qquad
 G_p=\gamma\Delta p-\gamma\Pi x+\frac{\gamma}{N}\mathbf1\otimes c.
\end{equation}
Positive row scaling preserves the product-set optimality conditions.
For \(u=x-x'\), \(v=p-p'\), strong convexity and completion of squares give
\begin{equation}\label{eq:monotone}
\begin{aligned}
 \langle G(s)-G(s'),s-s'\rangle
 &\ge m\|\Pi u\|^2+(m-\gamma\kappa^2/4)\|\Delta u\|^2\\
 &\quad+\gamma\|\Delta v-(\kappa/2)\Delta u\|^2.
\end{aligned}
\end{equation}
This holds for any corresponding subgradient selections. Condition
\eqref{eq:threshold} also implies \(m>\gamma/(N-1)^2\), so the own
allocation-price cost is jointly strongly convex: its price quadratic
has curvature \(\gamma\) and cross coefficient \(-\gamma/(N-1)\).
Thus own optimality conditions are sufficient.

% Ledger: 17, 19.
Substituting a KKT pair into~\eqref{eq:operator} supplies a game equilibrium.
Add its variational inequality to that of any other equilibrium.
The left side of~\eqref{eq:monotone} is then nonpositive, forcing
\(u=0\) and \(\Delta v=0\). At common prices the remaining game
conditions are precisely~\eqref{eq:kkt}. Strong convexity gives allocation
uniqueness.

% Ledger: 8, 9, 17, 19.
Equation~\eqref{eq:old-transfer}, the correction and complementarity yield
the transfer identity and budget balance. Choose
\(g_n\in\partial f_n(x_n)\), \(v_n\in N_{X_n}(x_n)\) with
\(g_n+\lambda+v_n=0\). Feasibility of zero gives \(v_n^\top x_n\ge0\).
The subgradient inequality at zero yields
\(0\ge f_n(x_n)+\lambda^\top x_n\), proving participation.
\end{proof}

\section{The fixed risk surrogate and economic errors}
% Ledger: 9, 13, 17, 19, 20, 22, 24.
Assume \(0\in X_n\). Losses \(J_n(x,\xi)\) are bounded in absolute
value by \(U_n\), samplewise \(L_n\)-Lipschitz and
\(m_n\)-strongly convex on \(X_n\). Samples are fresh and iid from a
query-independent distribution. For upper-tail probability
\(\beta_n\in(0,1)\), define
\begin{align}
 C_n(x)&=\inf_{\tau\in\mathbb R}\left\{\tau+
 \beta_n^{-1}\mathbb E[J_n(x,\xi)-\tau]_+\right\},\label{eq:cvar}\\
 \widehat C_{n,s_n}(x)&=\inf_{\tau\in\mathbb R}\left\{\tau+
 \frac1{\beta_ns_n}\sum_{i=1}^{s_n}[J_n(x,\xi_i)-\tau]_+\right\}.
\end{align}
Here \([a]_+=\max\{a,0\}\). True normalised valuation is
\(V_n(x)=C_n(0)-C_n(x)\). For satisfaction data the loss is
\(J_n=-\psi_n\); cash translation allows deterministic taxes outside the
risk functional. The planner sums individual risks, rather than taking
the risk of an aggregate random loss.

\subsection{One-sided value error and feasible queries}
% Ledger: 13, 17, 20, 22, 24, 26.
Suppose a known relative interior ball
\(a_n+\rho_n\mathcal B_n\subset X_n\) is available, where
\(\mathcal B_n\) is the unit ball in the local affine hull, of dimension
\(d_n>0\). Fix \(0<\varepsilon_n<1\) and
\(0<\delta_n<\varepsilon_n\rho_n\). Define
\begin{align}
 q_n(x,v)&=(1-\varepsilon_n)x+\varepsilon_na_n+\delta_nv,\\
 K_n(x)&=\mathbb E_{v,\xi^{1:s_n}}\widehat C_{n,s_n}(q_n(x,v)),
 \qquad v\sim\operatorname{Unif}(\mathcal B_n).\label{eq:K}
\end{align}
Decisions remain in \(X_n\), including zero. Set
\begin{equation}\label{eq:error-constants}
 e_n=\frac{U_n}{\beta_n}\sqrt{\frac{2\pi}{s_n}},\quad
 R_n=\sup_{x\in X_n}\|x-a_n\|,\quad
 w_n=e_n+2L_n\varepsilon_nR_n+L_n\delta_n.
\end{equation}
% Ledger: 13, 17, 20, 22, 24.
\begin{lemma}\label{lem:oracle}
Every query in~\eqref{eq:K} is locally feasible, and
\begin{equation}\label{eq:interval}
 -e_n-L_n\varepsilon_nR_n\le K_n(x)-C_n(x)
 \le L_n\delta_n+L_n\varepsilon_nR_n.
\end{equation}
The surrogate has curvature \(m_n(1-\varepsilon_n)^2\). With a uniform
sphere direction \(u\), independent of the fresh batch, the oracle
\begin{equation}\label{eq:oracle}
 \widehat g_n(x)=\frac{(1-\varepsilon_n)d_n}{\delta_n}
 \widehat C_{n,s_n}(q_n(x,u))u
\end{equation}
is unbiased for \(\nabla K_n(x)\), the relative gradient. It satisfies
\[
 \|\widehat g_n(x)\|\le (1-\varepsilon_n)d_nU_n/\delta_n.
\]
\end{lemma}
\begin{proof}
% Ledger: 20, 22, 24.
The query is the convex combination
\((1-\varepsilon_n)x+\varepsilon_n(a_n+(\delta_n/\varepsilon_n)v)\).
The strict radius inequality also leaves a relative neighbourhood on which
the smoothing identity is valid. Applying P's sphere identity and the chain
rule gives~\eqref{eq:oracle}. Since \(0\in X_n\), the affine hull is a
linear subspace; gradients can be lifted by an orthonormal basis.

% Ledger: 13, 17; source: P appendix, CVaR error calculation.
Let \(H_n=\mathbb E\widehat C_{n,s_n}\). Interchanging infimum and
expectation in the inequality direction gives \(H_n\le C_n\).
P's integrated Dvoretzky--Kiefer--Wolfowitz calculation gives
\(\mathbb E|\widehat C_{n,s_n}(x)-C_n(x)|\le e_n\), hence
\(C_n-e_n\le H_n\le C_n\).
At a centre \(y\) admitting ball queries, Jensen's inequality and
Lipschitz continuity therefore give
\begin{equation}\label{eq:centred}
 -e_n\le\mathbb E_v H_n(y+\delta_nv)-C_n(y)\le L_n\delta_n.
\end{equation}
Taking \(y=(1-\varepsilon_n)x+\varepsilon_na_n\) proves
\eqref{eq:interval}. These compare deterministic expected functions;
no expected supremum of random empirical errors is used.

% Ledger: 9, 17, 19, 20, 22, 24.
Samplewise strong convexity passes through CVaR by monotonicity,
convexity and cash translation: the deterministic quadratic deficit
passes through the risk functional unchanged. This applies also to
empirical CVaR and expectation. Affine contraction scales the deficit
by \((1-\varepsilon_n)^2\).
\end{proof}

\subsection{Economic consequence}
% Ledger: 13, 17, 19, 24, 26.
For a real function \(a\), let \(\osc_X a=\sup_Xa-\inf_Xa\).
Equation~\eqref{eq:interval} gives \(\osc_{X_n}(K_n-C_n)\le w_n\).
Use surrogate valuation \(V_n^K(x)=K_n(0)-K_n(x)\), and retain
payment~\eqref{eq:repair}.
% Ledger: 13, 17, 19, 24, 26.
\begin{theorem}\label{thm:economic}
Assume a surrogate planner KKT pair exists and
\[
 m_{\rm eff}=\min_n m_n(1-\varepsilon_n)^2
 >\frac{\gamma(2N-1)^2}{4(N-1)^4}.
\]
At every surrogate equilibrium \(\widetilde s\), its allocation
\(\widetilde x\) uniquely minimises \(\sum_nK_n\) on \(\mathcal A\).
Budget balance and surrogate individual rationality are exact. Writing
\(U_n(s)=V_n(x_n)-t_n^{\rm corr}(s)\), population-risk guarantees are
\begin{align}
 \sup_{s_n'\in X_n\times\mathbb R_+^r}
 \{U_n(s_n',\widetilde s_{-n})-U_n(\widetilde s)\}&\le w_n,\label{eq:regret}\\
 U_n(\widetilde s)&\ge-w_n,\label{eq:ir}\\
 \sum_n C_n(\widetilde x_n)-\sum_n C_n(x_n^\star)&\le\sum_nw_n,
 \label{eq:welfare}
\end{align}
where \(x^\star\) minimises population cost on \(\mathcal A\).
If total population cost has curvature \(m_0>0\), then
\begin{equation}\label{eq:allocation}
 \|\widetilde x-x^\star\|^2\le\frac{2}{m_0}\sum_nw_n.
\end{equation}
\end{theorem}
\begin{proof}
% Ledger: 13, 17, 24.
Lemma~\ref{lem:oracle} and Theorem~\ref{thm:implementation} give exact
surrogate implementation and transfers. Put \(a_n=K_n-C_n\) in this
proof only. The true minus surrogate payoff at allocation \(x_n\) is
\(a_n(x_n)-a_n(0)\). Its change under any unilateral deviation is at
most \(\osc a_n\), proving~\eqref{eq:regret}; its absolute value is at
most \(w_n\), proving~\eqref{eq:ir}. Transfers themselves do not change
when evaluated under true preferences, so budget balance remains exact.
Surrogate optimality gives
\[
 \sum_n[C_n(\widetilde x_n)-C_n(x_n^\star)]
 \le\sum_n[a_n(x_n^\star)-a_n(\widetilde x_n)]\le\sum_nw_n.
\]
Strong convexity and constrained optimality at \(x^\star\) lower-bound
this gap by \((m_0/2)\|\widetilde x-x^\star\|^2\).
\end{proof}
% Ledger: 13, 17, 24, 26.
Taking \(\delta_n=\varepsilon_n\rho_n/2\) yields error order
\(O(s_n^{-1/2}+\varepsilon_n)\), for fixed problem constants.
Where centred queries are already permissible,~\eqref{eq:centred} gives
the smaller oscillation bound \(e_n+L_n\delta_n\).
The improvement is at the objective level and does not remove the
inverse-radius dependence of oracle variance.

\section{A replacement learner}
% Ledger: 11, 17, 19, 24, 26; prior method: Cui--Shanbhag, Section 3.2.
Fix all surrogate parameters. Assume a compact price box
\([0,P_{\max}]^r\) containing a surrogate KKT equilibrium, and put
\(\mathcal S=\prod_n(X_n\times[0,P_{\max}]^r)\).
Comparison with that equilibrium in~\eqref{eq:monotone} shows every
boxed equilibrium has its allocation and common prices. Capacity is
therefore feasible already; slack resources force zero price by the
boxed price condition. No strict interiority of the reference price is
needed. This retains the unboxed equilibrium characterisation.

% Ledger: 11, 17, 19, 24.
Let \(G^K\) denote~\eqref{eq:operator} for the smooth costs
\(K_n-K_n(0)\). For \(\eta>0\), define
\[
 R(z)=(I+\eta(G^K+N_{\mathcal S}))^{-1}z,
 \qquad H_z(s)=G^K(s)+\eta^{-1}(s-z).
\]
The continuous monotone operator on the compact convex product gives a
well-defined resolvent. To approximate it, use projected stochastic
updates with the oracle~\eqref{eq:oracle}, analytic payment gradients,
and aggregate allocation and price messages:
\[
 s_{t+1}=\operatorname{proj}_{\mathcal S}
 (s_t-\alpha_t\widehat H_z(s_t)),\qquad
 \alpha_t=\frac{1}{\sigma(t+2)},\quad \sigma=\eta^{-1}.
\]
% Ledger: 11, 17, 19, 24.
\begin{proposition}\label{prop:learning}
Run \(T_k=\lceil(k+1)^{2+\zeta}\rceil\) inner updates centred at \(z_k\),
where \(\zeta>0\), and set \(z_{k+1}=s_{T_k}\). Then \(z_k\) converges
almost surely and in mean square to a possibly random surrogate equilibrium.
Its allocation converges in mean square to the deterministic optimum
\(\widetilde x\) of Theorem~\ref{thm:economic}.
\end{proposition}
\begin{proof}
% Ledger: 17, 19, 24.
Compactness and the oracle bound give a uniform conditional second moment
\(M^2\) for \(\widehat H_z\). With \(y=R(z)\), projection, the VI
at \(y\), and \(\sigma\)-strong monotonicity give
\[
 \mathbb E[\|s_{t+1}-y\|^2\mid s_t,z]
 \le(1-2\sigma\alpha_t)\|s_t-y\|^2+\alpha_t^2M^2.
\]
Induction yields
\[
 \mathbb E[\|s_T-R(z)\|^2\mid z]\le C/(T+1),\qquad
 C=\max\{\operatorname{diam}(\mathcal S)^2,M^2/\sigma^2\}.
\]
Thus the errors \(e_k=z_{k+1}-R(z_k)\) have
\(\sum_k\mathbb E\|e_k\|<\infty\), and are almost surely summable.

% Ledger: 17, 19, 24.
Subtracting the resolvent VIs proves firm nonexpansiveness. For any
fixed point \(z^\star\),
\[
 \|R(z_k)-z^\star\|^2\le\|z_k-z^\star\|^2-\|z_k-R(z_k)\|^2.
\]
Summable errors imply convergent distances to fixed points and vanishing
residual. Compactness supplies a cluster point; continuity places it in
the fixed-point set. The distance argument then gives convergence of the
whole sequence. Boundedness yields mean-square convergence by dominated
convergence. The allocation component is unique by
Theorem~\ref{thm:implementation}.
\end{proof}

\section{Related work and interpretation}
% Ledger: 12, 15, 19, 26; literature verification and reading: search.md.
The sampled CVaR one-point step is already published. Wang, Shen and
Zavlanos~\cite{wang2022}, Algorithm 1 and Lemmas 2--3, give the estimator,
smoothing comparison and distribution-function error bound, with regret
guarantees. The nested learner is also an application of an existing
principle: Cui and Shanbhag~\cite{cui2021}, Section 3.2, Propositions 5--6,
use stochastic inner solves of proximal inclusions and a polynomial
inner schedule of exponent greater than two to obtain almost-sure
convergence under monotonicity. We include the compact-case proof to
specify precisely which learner is justified here.

% Ledger: 15, 19, 26; literature comparison checked in search.md.
Wang et al.~\cite{wang2024}, Algorithm 1 and Theorem 1, use loss gradients
and estimated quantiles under strong monotonicity. Their theorem bounds
the time average of squared distance with high probability. Our weighted
operator retains a common-price degeneracy. Cherukuri~\cite{cherukuri},
Section 3, studies sample-average VIs whose operator components are
CVaRs, proving consistency and exponential solution-set approximation.
Here the implemented objective is the deterministic expectation of a
fixed empirical risk estimate, and the operator differentiates that
objective. Those statistical targets and guarantees differ.

% Ledger: 12, 13, 15, 17, 19, 24, 26.
The comparison supports a focused contribution: a correction of Q's
specific payment, with P's fixed surrogate interpreted through explicit
welfare and incentive errors. It does not support novelty claims for
CVaR game learning, the proximal construction or feasible smoothing in
general. The sources examined did not identify the same payment repair
and economic-bound package; this bounded search is not a proof of priority.
Economically, finite sampling changes the implemented preferences.
Equations~\eqref{eq:regret}--\eqref{eq:welfare} concern a fixed surrogate
equilibrium assessed using population risk, rather than a finite-time
regret guarantee for a learning trajectory.

\section{Limitations and open questions}
% Ledger: 9, 14, 17, 19, 24, 26.
The added assumptions are substantive: samplewise curvature, valid KKT
sum rules, existence of a surrogate planner multiplier, a feasible zero
outside action, known local interior balls and fresh losses at arbitrary
local queries. For learning, a known price cap must contain a KKT
price. Q's expected curvature assumption alone does not imply CVaR
curvature. Exact budget balance is asserted only at equilibrium; local
query feasibility does not ensure joint capacity feasibility of experiments.
No efficient total oracle complexity or stochastic benchmark is supplied.

% Ledger: 10, 14, 17, 19.
Weighted monotonicity does not prove Q's proximal best-response map
nonexpansive. In the recorded interior quadratic example with \(N=3\)
and \(m=\gamma=1\), its consensus-price column has squared norm
\(1+1/(2D^2)>1\), where \(D=(\mu+1)^2-1/4\), for every \(\mu>0\).
This defeats that certificate, without proving divergence of all damped
iterations. Proposition~\ref{prop:learning} concerns a changed algorithm.

% Ledger: 14, 17.
Allocation uniqueness likewise does not imply price uniqueness. For
\(N=3\), \(X_n=[0,1]\), \(c=3/2\), \(\gamma=1\), equiprobable
\(Z=\pm1\), and tail probability \(1/2\), the losses
\[
 J_n(x,Z)=x^2-3x+Z(x-1/2)
\]
give \(C_n(x)=x^2-3x+|x-1/2|\). The unique allocation is \(x_n=1/2\),
but \(\partial C_n(1/2)=[-3,-1]\) admits every common price in \([1,3]\).

% Ledger: 19, 24, 26.
Open questions are convergence of the original VS-PBR iteration, efficient
sample complexity for the replacement learner, a population-risk limit
with changing batches and radii, and jointly capacity-feasible exploration.
The fixed-parameter theorem establishes none of these. Broader publication
priority of the specific correction remains subject to further comparison.

\bibliographystyle{plain}
\bibliography{references}
\end{document}
